\scanpage{55}
\mspage{II 3}
% lines 1-9
En effet, \deleted{\unsure{cette} \unsure{représentation} \unsure{passe à un sous-groupe assez petit}} \inserted{l'hypothèse \inserted{\unsure{indique|implique} que} il existe un ss-groupe} ouvert $\pi$ de $\pi_1(X,\unsure{\xi|\zeta})$ tel que la représentation induite sur $I_{\unsure{\pi|1}}=\pi\cap\pi_1(\overline X,\xi)$ soit abélienne. Prenant $\unsure{U|\pi}$ assez petit, on peut supposer (I \quad) que \deleted{elle se \unsure{factorise par}} l'image de \deleted{\unsure{$\pi$}} dans $\operatorname{Aut}_{\mathbb{Q}_\ell}(\unsure{E|\mathbb{E}})$ est un pro-$\ell$-groupe. Donc la représentation \inserted{de $\pi$} provient d'une représentation de $\pi'$, \deleted{\unsure{extension}} l'extension (1.3) de $\pi_0$ par $I(\ell)$. On conclut par I 6.1.

% lines 10-18
\uline{Corollaire} 1.4. Avec les notations de 1.3, supposons l'action de $\pi_1(X,\xi)$ lui-même \deleted{\unsure{semi}} \inserted{essentiellement} abélienne. Alors \inserted{la représentation induite sur $I$ \deleted{\unsure{alors}} est essentiellement triviale,} \deleted{\unsure{En particulier}} \inserted{i.e.} il existe un \deleted{\unsure{revêtement galoisien} fini surjectif} \inserted{sous-groupe \deleted{distingué} ouvert $U$} \deleted{\unsure{connexe} $X'\mathrel{\text{---}}X$,} de $\pi$ tel que la restriction de la représentation envisagée à $\unsure{\pi|U}$ \deleted{\unsure{par}} se factorise en une représentation de \deleted{\unsure{$K$}} $\pi_0=\operatorname{Im}(\unsure{\pi_q|\pi}\mathrel{\text{---}}\pi_1(S,\xi))$.

% lines 19-29
\deleted{Cela résulte du fait que :}

On prend \deleted{\unsure{$U$}} d'abord $\pi$ tel que la restriction à $\pi$ soit abélienne, \deleted{donc par} et que \deleted{\unsure{provenant de}} l'image de $\pi$ dans $\operatorname{Aut}_{\mathbb{Q}_\ell}(E)$ soit un pro-$\ell$-groupe. Alors la représentation provient d'une représentation de $\pi^{\mathrm{ab}}(\ell)$, et \deleted{il \inserted{donc} suffit d'appliquer} 1.2 nous assure qu'il existe un sous-groupe ouvert $U$ de $\pi^{\mathrm{ab}}(\ell)$ tel que $U\mathrel{\text{---}}\pi_0^{\mathrm{ab}}(\ell)$ soit injectif, donc la représentation restreinte à $U$ provient d'une représentation de son image dans $\pi_0^{\mathrm{ab}}(\ell)$. Il suffit
% archive sheet number: 54
